% AGENT 2027 short BENCHMARK paper. DRAFT ONLY.
% Max 5 pages INCLUDING references, single-anonymous, IEEE conference style.
% Numbers must be regenerated/audited against results/summary_v3.json and RESULTS.md.
\documentclass[10pt,conference]{IEEEtran}
\usepackage{booktabs}
\usepackage{amsmath}
\usepackage{url}
\usepackage{graphicx}
\usepackage{balance}
\usepackage{float}
\begin{document}

\title{TrustVal: An Engineering Case Study of Agent Recovery After Ambiguous Writes}
\author{\IEEEauthorblockN{Vedant D. Borkute}
\IEEEauthorblockA{University of California, Riverside}}
\maketitle

\begin{abstract}
An API can commit a write even when its caller receives a timeout. How an agent responds depends partly on what the evaluation exposes and verifies. We report TrustVal, a four-task executable benchmark with a separate state service and a final-database verifier. In its first version, error messages revealed commit outcomes and server-side deduplication hid unsafe retries; two evaluated agents each passed all four tasks. After a documented redesign, Claude Haiku 4.5 and Sonnet 4.6 passed 19/40 and 35/40 task-trials, respectively, over ten runs of the same four tasks. All 26 failed trials violated exactly-once constraints. Saved traces show both blind retries and recovery through transaction-history checks. Newer evaluated agents nearly saturated the revised benchmark, including a delayed-visibility variant. The redesign changed several factors at once, and the injected transport errors have no matched uncommitted counterparts. TrustVal therefore documents how benchmark contracts can mask or reveal unsafe recovery decisions, rather than establishing general transaction-recovery ability.

\end{abstract}

\section{Introduction}
In a saved TrustVal run, an agent issued a supplier payout, received HTTP 504, and issued the same payout again. Both writes committed. A second agent encountering the same fault checked the payout record before proceeding and avoided the duplicate. Both could describe their actions in a plausible final response; only the durable state established whether the task was correct.

The distinction matters when testing tool-using agents. A non-idempotent API may apply a request even if the acknowledgement fails. Retrying then risks a second effect, but treating every timeout as a successful commit risks leaving work unfinished. Prior work already studies this problem across more services, faults, and harnesses \cite{limbo2026,undobench2026}. The question here is narrower: whether a compact executable evaluation actually exposes the behavior it purports to grade.

TrustVal began with two features that obscured recovery decisions. Its error text revealed whether a write had committed, and the server suppressed duplicates. After observing saturated scores, we revised the fault contract and grading setup, then examined saved database outcomes and tool traces.

We ask \textbf{RQ1}: Which implementation choices made the first test unable to distinguish unsafe retrying? \textbf{RQ2}: What did the revised test expose about the evaluated agents, and where did it cease to discriminate? We contribute the runnable four-task testbed, a documented instrument revision, and an evidence-backed analysis of the resulting failures and ceiling effects. This is a benchmark-engineering case study, not a new exactly-once protocol.

\section{Benchmark Design}
\subsection{Execution and grading}
Each task gives a terminal agent a command-line client for a small HTTP API. The agent container cannot inspect the API's SQLite database directly; a separate service owns that database. After the run, a hidden verifier reads the saved state and checks the requested records, absence of duplicates, and preservation of unrelated state. The agent's final message does not determine success.

Four tasks represent different observable states: a payout committed before HTTP 504; a role grant followed by a stale membership read, with a separate authoritative status path; two of three batch reservations committed before HTTP 504; and an account created before HTTP 502. An injected failure occurs once at a fixed point in each affected task. The resulting score means that these synthetic business-state invariants hold, not that the agent guarantees exactly-once execution in arbitrary distributed systems.

\subsection{Instrument revisions and controls}
Version 2 (v2) included outcome-revealing error text and server-side deduplication. Haiku 4.5 and Sonnet 4.6 each passed 4/4 tasks in one run. A correct final state thus did not establish that the agent had made a safe recovery decision: a blind retry could be harmless.

For version 3 (v3), we replaced the error text with ordinary HTTP 502/504 responses, removed server-side duplicate suppression and task-specific recovery hints, and exposed authoritative status and transaction-history reads. Client-provided references identify records for lookup but do not make writes idempotent. The v3 protocol was documented before the counted model runs. A scripted blind-retry control failed 4/4 tasks; a state-checking reference and the hosted oracle each passed 4/4.

The reference for the reservation task uses single-item writes rather than the fault-triggering batch endpoint, which the measured agents used. It validates an available safe path and the final-state verifier, not the reference's recovery from a failed batch. Several properties changed between v2 and v3, so the score difference cannot be assigned to any one change. Version 4 (v4) later added delayed read visibility. Table~\ref{tab:versions} summarizes these revisions.

% Version contrasts are descriptive. Several factors changed together between v2/v3.
\begin{table}[H]
\centering
\caption{How the four-task test changed. A control is a fixed program used to check the test; a reference agent is a prewritten solution.}
\label{tab:versions}
\scriptsize
\renewcommand{\arraystretch}{1.10}
\begin{tabular}{@{}p{0.11\columnwidth}p{0.40\columnwidth}p{0.37\columnwidth}@{}}
\toprule
Version & Ambiguous outcome and tool behavior & Validation and observed limit \\
\midrule
v2 & Error messages revealed whether a change was saved; the service ignored repeated requests. & Haiku 4.5 and Sonnet 4.6 each passed 4/4 in one run. \\
v3 & Ordinary HTTP 502/504 errors; repeated writes could create duplicates; status and saved-record checks available; no recovery hints. & Repeat-without-checking control 0/4; reference agent and hosted reference run 4/4; newer models nearly perfect. \\
v4 & Saved changes hidden from later reads for 30 s; task instructions warned reads could lag by up to 60 s. & Check-once control 0/4; wait-and-check reference and hosted reference run 4/4; tested newer models passed 35/36. \\
\bottomrule
\end{tabular}
\end{table}


\section{Experimental Protocol}
We ran the tasks in Valkyrie/Harbor with the Modal VM provider. Claude models used the official Claude Code harness, and OpenAI models used the Codex harness. The GPT-5.x dataset has an additional OpenAI API network allowlist; otherwise the v3 tasks and verifiers match the Claude version. Cross-provider results therefore also reflect different harnesses and cannot isolate model effects.

Seven models were each evaluated in ten separate runs of the same four fixed tasks, one trial per task per run. Three GPT-6-family models were tested once each as exploratory checks and are excluded from the main table and model comparison. The outcome is a binary hidden-database verifier score. Saved databases were independently regraded and matched the hosted grades. We distinguish agent errors from setup failures and oracle/control runs. Run-level scores, not the four correlated tasks within a run, are the units used in the reported permutation comparison.

\section{Results}
% Generated from results/summary_v3.json. Do not hand-edit counts.
% CT=timeout after payout was saved; SR=old membership returned after role change; PB=two of three reservations saved before timeout; DD=repeat request saved twice.
\begin{table}[H]
\centering
\caption{Passes on the same four v3 tasks, each repeated ten times per model. A pass means the saved records met the task rules. CT: timeout after a payout was saved; SR: old membership information returned after a role change; PB: two of three reservations saved before timeout; DD: a repeated request saved a second copy. Different providers used different agent software to run and record the models.}
\label{tab:v3}
\scriptsize
\setlength{\tabcolsep}{2.5pt}
\begin{tabular}{@{}lrrrrr@{}}
\toprule
Model & CT & SR & PB & DD & Total \\
\midrule
Claude Haiku 4.5 & 6/10 & 10/10 & 3/10 & 0/10 & 19/40 \\
Claude Sonnet 4.6 & 10/10 & 10/10 & 7/10 & 8/10 & 35/40 \\
Claude Haiku 5.5 & 10/10 & 10/10 & 10/10 & 9/10 & 39/40 \\
Claude Sonnet 5.5 & 10/10 & 10/10 & 10/10 & 10/10 & 40/40 \\
Claude Opus 5.5 & 10/10 & 10/10 & 10/10 & 10/10 & 40/40 \\
GPT-5.3-Codex & 10/10 & 10/10 & 9/10 & 5/10 & 34/40 \\
GPT-5.4 & 10/10 & 10/10 & 10/10 & 10/10 & 40/40 \\
\bottomrule
\end{tabular}
\end{table}

\subsection{Duplicate effects in v3}
Table~\ref{tab:v3} gives task-level passes across ten runs of four fixed scenarios per model. Haiku 4.5 passed 19/40 and Sonnet 4.6 passed 35/40; the average difference was 1.6 tasks per four-task run (two-sided exact run-level permutation $p=2.3\times10^{-4}$). The calculation contrasts these observed runs, not a population of unseen workflows.

Across the two Claude 4.x models, all 26 failed trials had \texttt{exactly\_once=0} in the saved regrades. The run records classify these failures as blind retries following committed writes; we independently checked representative tool traces, not every trajectory. Haiku 4.5 acknowledged duplication in 17 of its 21 failures. Sonnet 4.6 claimed full success in four of five failures. Final-response confidence and final-state correctness can therefore disagree.

GPT-5.3-Codex passed 34/40, with duplicates in all six failed tasks. Haiku 5.5 passed 39/40; its single failure was an argument-order mistake without an injected fault. Sonnet 5.5, Opus 5.5, and GPT-5.4 each passed 40/40. Three GPT-6-family pilots each passed 4/4 in one run and are excluded from Table~\ref{tab:v3}. These are fixed-task observations, not a ranking over unseen tasks or matched model-and-harness conditions.

\subsection{Two ways to respond to HTTP 504}
In Haiku 4.5 run \texttt{57fcbdfa}, the first payout for invoice INV-104 committed but returned HTTP 504. The agent paid INV-105, then retried INV-104 without checking its status. The database contained two separate committed operations for INV-104, and the agent subsequently acknowledged the overpayment. In Sonnet 4.6 run \texttt{7f183c15}, the agent checked the transaction history and reference status after the same fault, confirmed the INV-104 payment had already committed, and paid only INV-105. This run passed. Table~\ref{tab:trace} shows the contrast.

% Verbatim run IDs and observed outcomes; see TRACE_AUDIT.md for provenance.
\begin{table}[t]
\centering
\caption{Different recovery actions after HTTP 504 on payout INV-104 in v3.}
\label{tab:trace}
\scriptsize
\setlength{\tabcolsep}{3pt}
\renewcommand{\arraystretch}{1.07}
\begin{tabular}{@{}p{0.23\columnwidth}p{0.42\columnwidth}p{0.23\columnwidth}@{}}
\toprule
Model and run & After the timeout & Final ledger \\
\midrule
Haiku 4.5 \texttt{57fcbdfa} & Paid INV-105, then retried INV-104 without a state read. & Two INV-104 payouts; fail. \\
Sonnet 4.6 \texttt{7f183c15} & Read status and ledger, then paid only INV-105. & One per invoice; pass. \\
\bottomrule
\end{tabular}
\end{table}


False reporting is distinct from the duplicate itself. In Sonnet 4.6's partial-batch run \texttt{11aec2b6}, two reservations committed before HTTP 504. The agent retried the complete batch without reading the existing reservations, then reported that all three items had been committed exactly once. The database regrade failed the task. A completion-message-only evaluator would have missed the state violation.

\subsection{Delayed visibility and a near-ceiling result}
In v4, committed writes remained invisible to reads for 30 seconds; the client documentation warned that visibility could take up to 60 seconds. A waiting reference and the hosted oracle passed all four tasks, whereas a check-once retry policy failed. Across three four-task runs each, Haiku 5.5 passed 11/12, Sonnet 5.5 passed 12/12, and GPT-6.1 Terra passed 12/12. The single failure was partly shaped by Claude Code's restrictions on waiting and background commands. These small samples do not support a between-model conclusion.

\section{Related Work and Positioning}
Stateful tool benchmarks, including $\tau$-bench \cite{taubench2025} and ToolSandbox \cite{toolsandbox2025}, grade outcomes beyond surface text. Closer studies directly address recovery from non-atomic effects. Verified Tool Calls \cite{verifiedcalls2026} evaluates verification-aware wrappers, while AFT-Bench \cite{aftbench2026} holds the backend and fault fixed to isolate interface interventions.

LIMBO \cite{limbo2026} spans six services, twelve fault modes, and 25,930 episodes, including delayed effects and contract-level idempotency. It also studies when read-back cannot establish exactly-once execution. UndoBench \cite{undobench2026} pairs nominal and faulted workflows with state and effect-history checks; Evidence-Guided Recovery \cite{evidenceguided2026} examines recovery choices under partial execution. AgentFaultBench \cite{agentfaultbench2026} evaluates fault recovery in a stateful GitLab environment.

TrustVal is smaller and does not isolate harness or interface effects as those studies do. Its intended value is an inspectable record of a benchmark whose initial implementation could not reveal unsafe retries, together with revised tasks and original execution traces. The failure mode itself is not novel.

\balance
\section{Threats to Validity and Engineering Lessons}
\textbf{Scope and missing counterfactual.} All injected v3 errors follow at least a partial commit; we did not pair each error with an observationally identical \emph{uncommitted} outcome. On the single-write tasks, simply never retrying after an error could succeed without establishing that an agent diagnosed what happened. The four scenarios also fix fault location and initial state; success cannot be generalized to unseen APIs, actual payments, or arbitrary retry policies. A full ambiguity benchmark would add matched committed and uncommitted branches, then condition recovery on nominal-task competence.

\textbf{Instrument and analysis.} We changed several v2 properties together after inspecting its saturated runs; v4 was designed after observing the v3 ceiling. Neither transition identifies a single causal factor. The v3 \texttt{stale\_read} scenario produced no transport error and was solved in every run. The batch reference solution avoids the fault-triggering batch command by using individual reservations, although measured agents used the batch endpoint. The tasks are public, and model familiarity cannot be ruled out. Ten runs repeat the same four tasks, not a sample of forty distinct workflows; the permutation statistic refers only to the observed run distribution. Claude and Codex used different harnesses, so cross-provider contrasts are not isolated model effects. The v4 three-run samples cannot resolve small differences, and one failure depended partly on harness wait behavior. No reversal operation was offered.

\textbf{Reusable checks.} A stateful agent benchmark should (1) distinguish the tool's response from durable effects; (2) verify that a naive policy fails when it should and that the reference actually exercises the fault; (3) preserve traces and grade post-run state independently of final prose; and (4) report saturated scenarios and protocol revisions alongside positive findings. TrustVal demonstrates the need for these checks, rather than proving a new recovery method.

\subsection{Reproducibility}
The public artifact at \url{https://github.com/Vedant2100/trustval} includes versioned task manifests, sidecar services, naive and reference policies, saved run identifiers, and regrade scripts. The scripts run verifiers against exported SQLite files and compare scores with Valkyrie's reports. Earlier infrastructure failures are recorded separately from model trials. Rerunning hosted agents requires container support, compatible Harbor/Valkyrie tooling, and model API access.

\section{Conclusion}
TrustVal's first version could pass an unsafe retry because the error disclosed the commit and the API suppressed duplicates. Removing those protections exposed duplicate writes in some measured agent runs, while newer agents nearly saturated the resulting four-task test. The artifact is a regression check and a concrete example of why benchmark validity must be tested with adversarial controls and independent state inspection. General recovery from ambiguous outcomes remains unmeasured without matched uncommitted cases.

\bibliographystyle{IEEEtran}
\bibliography{references}
\end{document}
